\documentclass[12pt]{article}
\usepackage[utf8]{inputenc}
\usepackage[T1]{fontenc}
\usepackage[margin=1in]{geometry}
\usepackage{setspace}
\usepackage{xcolor}
\usepackage[colorlinks=true, urlcolor=blue!60!black, linkcolor=blue!60!black]{hyperref}
\setlength{\parindent}{2em}
\setlength{\parskip}{1.25em}
\renewcommand{\baselinestretch}{1.0}
\title{\textit{Penn Central Transportation Co. v. New York City}, 438 U.S. 104 (1978)}
\author{Case Brief}
\date{}
\begin{document}
\maketitle

\section{Case Metadata}
\begin{itemize}
  \item \textbf{Citation:} \textit{Penn Central Transportation Co.\ v.\ New York City}, 438 U.S.\ 104 (1978) (No.\ 77-444).
  \item \textbf{Argued / Decided:} Argued April 17, 1978; Decided June 26, 1978.
  \item \textbf{Vote:} 6--3 (affirming the New York Court of Appeals).
  \item \textbf{Majority Author:} Justice William J.\ Brennan, Jr.\ (joined by Justices Stewart, White, Marshall, Blackmun, and Powell).
  \item \textbf{Concurrences:} None.
  \item \textbf{Dissent:} Justice William H.\ Rehnquist (joined by Chief Justice Burger and Justice Stevens).
  \item \textbf{Chapter / Topic:} Chapter 6: Substantive Due Process --- Takings \& Property Rights (Regulatory Takings, the Three-Factor Balancing Test, and the ``Parcel as a Whole'' Rule).
  \item \textbf{Primary Source:} \href{https://supreme.justia.com/cases/federal/us/438/104/}{Justia U.S.\ Supreme Court Center (438 U.S.\ 104)}; see also \href{https://www.law.cornell.edu/wex/takings}{Cornell LII Wex: Takings}.
\end{itemize}

\section{Facts and Procedural History}
In 1965, New York City enacted its Landmarks Preservation Law (N.Y.C.\ Admin.\ Code, ch.\ 8-A, \S\ 201.0 et seq.), empowering an 11-member Landmarks Preservation Commission (subject to Board of Estimate modification and judicial review) to designate historic districts and individual ``landmarks'' on ``landmark sites.'' Designation required owners to keep exteriors ``in good repair'' and obtain Commission approval---via a certificate of ``no exterior effect,'' ``appropriateness,'' or ``insufficient return''---before making exterior alterations. To soften the burden, New York City's zoning ordinances (liberalized in 1968 and 1969) granted landmark owners Transferable Development Rights (``TDRs'') allowing unused floor-area air rights to be transferred to proximate lots under common ownership.

On August 2, 1967, the Commission designated Grand Central Terminal---an eight-story, 1913 French Beaux-Arts railroad station in midtown Manhattan owned by Penn Central Transportation Co.\ and its affiliates---as a landmark, and its city tax block as a landmark site. Penn Central did not seek judicial review of the designation. In January 1968, to boost revenue, Penn Central entered a renewable 50-year lease with UGP Properties, Inc., under which UGP would construct a multistory office tower above the Terminal and pay Penn Central \$1 million annually during construction and at least \$3 million annually thereafter (partially offset by \$700,000 to \$1 million in displaced concession rents).

Penn Central and UGP submitted two zoning-compliant designs by architect Marcel Breuer: \textit{Breuer I} (a 55-story office tower cantilevered above the existing facade) and \textit{Breuer II Revised} (a 53-story tower that tore down part of the 42d Street facade). The Commission denied certificates of no exterior effect and appropriateness for both, stating that tearing down a facade destroys a landmark and calling the 55-story cantilevered tower an ``aesthetic joke'' that would overwhelm the Terminal's Beaux-Arts setting. Without seeking judicial review of the denials or proposing a smaller structure, Penn Central and UGP sued in New York Supreme Court, Trial Term, alleging an uncompensated taking under the Fifth and Fourteenth Amendments and a Fourteenth Amendment due process violation. The trial court granted declaratory and injunctive relief, but the Appellate Division reversed (50 App.\ Div.\ 2d 265 (1975)), and the New York Court of Appeals affirmed (42 N.Y.2d 324 (1977)), finding no taking because Penn Central could still earn a reasonable return and held valuable TDRs. The U.S.\ Supreme Court noted probable jurisdiction and affirmed.

\section{Issue(s) Presented}
\begin{enumerate}
  \item Whether New York City's application of its Landmarks Preservation Law to prohibit construction of a 50+ story office tower atop Grand Central Terminal effects a regulatory ``taking'' of private property under the Fifth and Fourteenth Amendments.
  \item Whether a regulatory takings claim may be evaluated by severing a single property interest---here, superjacent air rights---from the owner's ``parcel as a whole.''
\end{enumerate}

\section{Holding}
No to both issues (6--3). Applying New York City's Landmarks Preservation Law to Grand Central Terminal did not effect a Fifth and Fourteenth Amendment taking. Establishing the default framework for regulatory takings, the Court held that courts must conduct an ad hoc, factual inquiry weighing three primary factors---(1) the economic impact of the regulation on the claimant, (2) the extent of interference with distinct investment-backed expectations, and (3) the character of the governmental action---measured against the ``parcel as a whole'' rather than conceptually severed air rights. Because the law preserved the Terminal's profitable existing use, allowed Penn Central a reasonable return, did not bar a smaller harmonizing addition, and conferred valuable TDRs, no compensation was required.

\section{Reasoning}
Writing for the Court, Justice Brennan observed that while the Takings Clause bars government from forcing a few people alone to bear public burdens that in fairness and justice should be borne by the public as a whole (\textit{Armstrong v.\ United States}, 364 U.S.\ 40, 49 (1960)), the Court has developed no ``set formula'' for when a regulation ``goes too far'' (\textit{Pennsylvania Coal Co.\ v.\ Mahon}, 260 U.S.\ 393, 415 (1922)). Instead, Brennan synthesized prior land-use cases into the three-factor ad hoc inquiry---economic impact, interference with distinct investment-backed expectations, and character of the governmental action---noting that a taking is far more readily found when government causes a physical invasion (\textit{United States v.\ Causby}, 328 U.S.\ 256 (1946)) than when a public program adjusts economic benefits and burdens for the common good.

Brennan then rejected four broad theories advanced by Penn Central:
\begin{enumerate}
  \item \textbf{Rejection of Conceptual Severance (``Parcel as a Whole''):} Penn Central could not manufacture a total taking by isolating its superjacent air rights. Takings analysis focuses on the character and extent of interference with the entire city tax block designated as the landmark site, not discrete vertical or lateral segments.
  \item \textbf{Not ``Reverse Spot'' Zoning:} Unlike arbitrary spot zoning, the Landmarks Law implemented a comprehensive citywide preservation plan covering over 400 landmarks and 31 historic districts, backed by judicial review.
  \item \textbf{Disproportionate Burden Does Not Establish a Taking:} General welfare laws (such as zoning in \textit{Euclid v.\ Ambler Realty Co.} or nuisance-prevention in \textit{Hadacheck v.\ Sebastian}) routinely burden some owners more severely than others. Moreover, Penn Central shared in the citywide economic and cultural benefits of landmark preservation.
  \item \textbf{No Enterprise Appropriation:} Unlike military overflights in \textit{Causby}, the City did not appropriate the airspace for governmental operations; it merely restricted private alteration of the exterior.
\end{enumerate}

Applying the three factors to the Terminal site, Brennan found the severity of impact insufficient to constitute a taking: (a) the law did not interfere with Penn Central's primary, 65-year use as a railroad terminal with offices and concessions, which remained capable of earning a ``reasonable return''; (b) the Commission had not foreclosed all airspace development, as Penn Central never proposed a smaller, harmonizing structure (such as the 20-story tower originally contemplated in 1913); and (c) even if TDRs might not alone constitute ``just compensation'' for a taking, their transferability to at least eight nearby Penn Central properties significantly mitigated the financial burden.

\section{Concurring and Dissenting Opinions}
\paragraph{Dissenting Opinion of Justice Rehnquist (joined by Chief Justice Burger and Justice Stevens).}
Justice Rehnquist argued that New York City's landmark designation effected a compensable taking. Distinguishing individual landmarking from traditional zoning or historic-district ordinances, Rehnquist emphasized that zoning creates ``an average reciprocity of advantage'' (\textit{Mahon}, 260 U.S.\ at 415) by mutually restricting and benefiting all neighbors, whereas New York singled out fewer than 0.1\% of its over one million buildings, imposing multimillion-dollar losses and an affirmative duty of upkeep backed by criminal penalties on a handful of owners for the public's benefit. Because Penn Central's proposed tower was fully zoning-compliant and no nuisance, the City was forcing Penn Central to bear a public burden simply because ``too good a job was done in designing and building'' the Terminal. Rehnquist would have held that a taking occurred and remanded to determine whether Penn Central's TDRs provided the ``full and perfect equivalent'' required for just compensation.

\section{Key Quotations}
\begin{enumerate}
  \item \textbf{The Three-Factor Balancing Test (Brennan, J., 438 U.S.\ at 124):}
  \begin{quote}
  ``In engaging in these essentially ad hoc, factual inquiries, the Court's decisions have identified several factors that have particular significance. The economic impact of the regulation on the claimant and, particularly, the extent to which the regulation has interfered with distinct investment-backed expectations are, of course, relevant considerations. \dots\ So, too, is the character of the governmental action. A `taking' may more readily be found when the interference with property can be characterized as a physical invasion by government \dots\ than when interference arises from some public program adjusting the benefits and burdens of economic life to promote the common good.''
  \end{quote}
  \item \textbf{The ``Parcel as a Whole'' Rule (Brennan, J., 438 U.S.\ at 130--31):}
  \begin{quote}
  ```Taking' jurisprudence does not divide a single parcel into discrete segments and attempt to determine whether rights in a particular segment have been entirely abrogated. In deciding whether a particular governmental action has effected a taking, this Court focuses rather both on the character of the action and on the nature and extent of the interference with rights in the parcel as a whole---here, the city tax block designated as the `landmark site.'''
  \end{quote}
  \item \textbf{Singling Out Few Owners for Public Burdens (Rehnquist, J., dissenting, 438 U.S.\ at 147):}
  \begin{quote}
  ``Appellees have imposed a substantial cost on less than one one-tenth of one percent of the buildings in New York City for the general benefit of all its people. It is exactly this imposition of general costs on a few individuals at which the `taking' protection is directed.''
  \end{quote}
  \item \textbf{Penalty for Architectural Excellence (Rehnquist, J., dissenting, 438 U.S.\ at 146):}
  \begin{quote}
  ``Penn Central is prevented from further developing its property basically because too good a job was done in designing and building it. The city of New York, because of its unadorned admiration for the design, has decided that the owners of the building must preserve it unchanged for the benefit of sightseeing New Yorkers and tourists.''
  \end{quote}
\end{enumerate}

\section{Doctrinal Significance and Modern Legacy (Through 2026)}
\textit{Penn Central} remains the foundational default test for regulatory takings in 2026, yet the Supreme Court has steadily carved out categorical \textit{per se} exceptions that bypass \textit{Penn Central} balancing:
\begin{itemize}
  \item \textbf{Categorical Per Se Exceptions --- \textit{Loretto} and \textit{Lucas}:} In \href{https://supreme.justia.com/cases/federal/us/458/419/}{\textit{Loretto v.\ Teleprompter Manhattan CATV Corp.}, 458 U.S.\ 419 (1982)}, the Court held that any permanent physical occupation is a \textit{per se} taking regardless of minimal economic impact. In \href{https://supreme.justia.com/cases/federal/us/505/1003/}{\textit{Lucas v.\ South Carolina Coastal Council}, 505 U.S.\ 1003 (1992)}, the Court held that a regulation depriving an owner of \textit{all} economically beneficial use of land is a \textit{per se} taking unless justified by background principles of state property or nuisance law; anything less than a 100\% wipeout reverts to \textit{Penn Central}.
  \item \textbf{Refining the ``Denominator'' --- \textit{Murr v.\ Wisconsin} (2017):} In \href{https://supreme.justia.com/cases/federal/us/582/15-214/}{\textit{Murr v.\ Wisconsin}, 582 U.S.\ 383 (2017)} (\href{https://www.law.cornell.edu/supremecourt/text/15-214}{Cornell LII}), the Court reaffirmed \textit{Penn Central}'s ``parcel as a whole'' principle and adopted a multifactor test (treatment under state/local law, physical characteristics, and prospective value under the regulation) to decide when adjacent lots should be merged as a single denominator.
  \item \textbf{Expanding Per Se Physical Takings --- \textit{Cedar Point} (2021) and \textit{Tyler} (2023):} In \href{https://supreme.justia.com/cases/federal/us/594/20-107/}{\textit{Cedar Point Nursery v.\ Hassid}, 594 U.S.\ 139 (2021)}, the Court sharply narrowed \textit{Penn Central}'s domain, holding that a California regulation granting union organizers intermittent access to farmland (up to three hours a day, 120 days a year) appropriates a right to invade and constitutes a \textit{per se} physical taking rather than a use restriction governed by \textit{Penn Central}. Similarly, in \href{https://supreme.justia.com/cases/federal/us/598/22-166/}{\textit{Tyler v.\ Hennepin County}, 598 U.S.\ 631 (2023)}, the Court treated a county's retention of surplus equity from a tax foreclosure sale as a classic physical confiscation rather than a regulatory taking.
  \item \textbf{Unconstitutional Conditions and Legislative Exactions --- \textit{Sheetz} (2024--2026):} In \href{https://supreme.justia.com/cases/federal/us/601/22-1074/}{\textit{Sheetz v.\ County of El Dorado}, 601 U.S.\ 267 (2024)}, a unanimous Court held that legislatively enacted land-use permit exactions (such as statutory traffic-impact fees) are not exempt from the heightened \textit{Nollan}/\textit{Dolan} ``essential nexus'' and ``rough proportionality'' scrutiny, preventing legislatures from routing exactions into deferential \textit{Penn Central} review. Through 2026, \textit{Penn Central} continues to govern ordinary use restrictions (zoning, historic preservation, rent and environmental limits), where landowners rarely prevail unless they can reframe the restriction as a physical appropriation under \textit{Cedar Point}, a total wipeout under \textit{Lucas}, or an exaction under \textit{Sheetz}.
\end{itemize}

\section{Pedagogical Critiques and Exam Hypotheticals}
\paragraph{Pedagogical Critiques.}
\begin{enumerate}
  \item \textbf{Indeterminacy and Pro-Government Tilt:} Critics note that \textit{Penn Central} offers no mathematical weights among its three factors, making outcomes unpredictable while heavily favoring regulators whenever a parcel retains any viable economic use.
  \item \textbf{The ``Denominator Problem'' and TDR Ambiguity:} Defining the ``parcel as a whole'' is outcome-determinative (a 100\% loss of a severed segment vs.\ a partial loss of the whole), yet \textit{Penn Central} and \textit{Murr} leave the denominator fluid. Likewise, counting TDRs on the ``taking'' side of the ledger risks letting governments offset takings with speculative non-cash credits rather than cash ``just compensation.''
\end{enumerate}

\paragraph{Exam Hypothetical 1 (Baseline --- No Taking Under \textit{Penn Central}).}
The City of Harborview enacts a comprehensive Historic Waterfront Ordinance designating 180 Victorian warehouses, including Bayside Lofts (owned by Harbor Corp.\ since 1995), as architectural landmarks. Harbor Corp.\ wants to demolish the upper two floors and build a 30-story luxury condominium tower, which would triple the building's value. The Preservation Board denies the permit, though Bayside Lofts remains fully leased as commercial offices generating a steady 7\% annual return, and the ordinance allows Harbor Corp.\ to transfer unused floor-area ratios to three adjacent commercial lots it owns.
\begin{itemize}
  \item \textbf{Analysis:} This falls squarely under \textit{Penn Central}, not \textit{Lucas} or \textit{Cedar Point}, because there is no physical invasion and no total loss of value. Measuring the parcel as a whole (the entire Bayside Lofts lot), Harbor Corp.\ cannot sever its air rights. All three \textit{Penn Central} factors favor Harborview: (1) economic impact is non-fatal given the 7\% return and mitigating TDRs; (2) the ordinance preserves Harbor Corp.'s primary investment-backed expectation (commercial leasing since 1995); and (3) the character of the action is a comprehensive land-use program rather than a physical occupation. No taking occurred.
\end{itemize}

\paragraph{Exam Hypothetical 2 (``Opposite Side'' Contrast --- Per Se Taking / Exaction Outside \textit{Penn Central}).}
Modify Hypothetical 1 on the \textbf{opposite side} of the doctrinal line: Harborview amends its ordinance to require every designated waterfront landmark owner (a) to grant the general public a permanent daytime easement along a 15-foot-wide strip of its private dock for sightseeing, and (b) if the owner applies for a routine interior renovation permit, to pay a flat \$250,000 statutory ``Harborwide Public Art Fee'' set by the City Council with no connection to the renovation's impact.
\begin{itemize}
  \item \textbf{Analysis:} Neither requirement is shielded by \textit{Penn Central} balancing. First, mandating a public sightseeing easement appropriates the owner's right to exclude and constitutes a \textit{per se} physical taking under \textit{Cedar Point Nursery v.\ Hassid}, 594 U.S.\ 139 (2021) (and \textit{Loretto}), regardless of how profitable the rest of the building remains. Second, under \textit{Sheetz v.\ County of El Dorado}, 601 U.S.\ 267 (2024), the City Council's legislatively enacted \$250,000 permit fee is not immune from \textit{Nollan}/\textit{Dolan} review; because a flat public-art fee bears no essential nexus or rough proportionality to an interior renovation's impacts, it is an unconstitutional exaction.
\end{itemize}

\end{document}
